\subsection{Future work}
Four extensions follow directly from the limitations above. First, running the \texttt{full}
experiment profile (Section~\ref{sec:appendix-profiles}) would replace our point estimates with
their paper-scale counterparts and let us report confidence intervals on the crossover
$\sigma^*$ itself rather than on the bias and RMSE curves that locate it. Second, verifying the
Eyam-1666 digitization against Raggett (1982) directly would let E14 be reported as a validated
historical estimate rather than a pipeline demonstration. Third, since our model variants (E12)
remain low-dimensional and non-stiff, testing whether the crossover phenomenon's specific
numerical values -- as opposed to its qualitative existence -- carry over to stiffer or
higher-dimensional compartmental models (e.g.\ age-structured or spatially coupled SIR variants)
is open. Fourth, our synthetic-data experiments assume the fitted model is correctly specified;
extending the noise model to misspecified dynamics (an SIR fit to SEIR-generated data, for
instance) would let the solver-vs-noise crossover be studied alongside a third error source,
model-form error, rather than the two considered here.
